\chapter{Training Data Generation Engine}
\label{chap:generator}

\section{Algorithm 2: Agentic Data Synthesis}

Given taxonomies $\mathcal{T}$ and a target size $N$, the generator produces a synthetic
dataset optimized for diversity, complexity, and quality.

\begin{figure}[H]
\centering
\begin{tikzpicture}[
    node distance=0.6cm and 1.5cm,
    stepbox/.style={rectangle, draw, rounded corners, fill=sapgold!10,
                 minimum width=3cm, minimum height=0.7cm, font=\small,
                 text width=2.8cm, align=center},
    arr/.style={-{Stealth[length=2mm]}, thick},
]
    \node[stepbox] (s1) {1. Mix Sampling};
    \node[stepbox, below=of s1] (s2) {2. Meta-Prompt Generation};
    \node[stepbox, below=of s2] (s3) {3. Complexification ($c=0.5$)};
    \node[stepbox, below=of s3] (s4) {4. Example Generation};
    \node[stepbox, below=of s4] (s5) {5. Double-Critic Filtering};
    
    \draw[arr] (s1) -- (s2);
    \draw[arr] (s2) -- (s3);
    \draw[arr] (s3) -- (s4);
    \draw[arr] (s4) -- (s5);
\end{tikzpicture}
\caption{Data generation stages implementing Algorithm 2.}
\label{fig:generator-stages}
\end{figure}

\section{Step 1: Mix Sampling (Global Diversity)}

Nodes are sampled from taxonomies to create ``mixes''---combinations of factors:

\begin{lstlisting}[language=Python,caption={Mix sampling (simula\_data\_generator.py, lines 120--150)}]
def _sample_mix(
    self, taxonomies: list[Taxonomy], strategy: SamplingStrategy
) -> list[TaxonomyNode]:
    """Sample a mix of nodes from taxonomies."""
    mix = []
    for taxonomy in taxonomies:
        if strategy == SamplingStrategy.UNIFORM:
            # Sample uniformly from leaf nodes
            leaves = self._get_leaf_nodes(taxonomy.root)
            node = random.choice(leaves)
        elif strategy == SamplingStrategy.COVERAGE:
            # Prioritize under-sampled nodes
            node = self._get_least_sampled_node(taxonomy)
        mix.append(node)
    return mix
\end{lstlisting}

\section{Step 2: Meta-Prompt Generation (Local Diversity)}

Each mix is converted to a meta-prompt that guides example generation:

\begin{lstlisting}[language=Python,caption={Meta-prompt generation}]
async def _generate_meta_prompt(
    self, mix: list[TaxonomyNode], schema_context: str
) -> str:
    """Generate a meta-prompt from a taxonomy mix."""
    prompt = f"""
    Generate a natural language question that a user might ask,
    given these requirements:
    
    {self._format_mix(mix)}
    
    Schema context:
    {schema_context}
    
    The question should be realistic and answerable with SQL.
    """
    response = await self.llm.generate(prompt, temperature=0.9)
    return response.content
\end{lstlisting}

\section{Step 3: Complexification}

A fraction $c$ (default 0.5) of examples undergo complexification:

\begin{lstlisting}[language=Python,caption={Complexification prompt}]
SYSTEM: You increase the complexity of a Text-to-SQL example
        while maintaining correctness.
        
USER:   Original question: {question}
        Original SQL: {sql}
        
        Add complexity by:
        - Adding edge cases or constraints
        - Requiring additional joins
        - Adding subqueries or CTEs
        - Handling NULL values explicitly
        
        Output the complexified question and SQL.
\end{lstlisting}

\section{Step 4: Example Generation}

The meta-prompt is used to generate a complete training example:

\begin{lstlisting}[language=Python,caption={Example generation}]
async def _generate_example(
    self, meta_prompt: str, schema_context: str
) -> TrainingExample:
    """Generate a single training example."""
    prompt = f"""
    Generate a Text-to-SQL training example.
    
    Meta-prompt: {meta_prompt}
    Schema: {schema_context}
    
    Output JSON:
    {{
      "question": "...",
      "sql": "SELECT ..."
    }}
    """
    response = await self.llm.generate(prompt)
    return self._parse_example(response.content)
\end{lstlisting}

\section{Step 5: Double-Critic Filtering}

The double-critic mitigates sycophancy bias by independently querying correctness
and incorrectness:

\begin{lstlisting}[language=Python,caption={Double-critic evaluation (simula\_llm\_client.py, lines 180--220)}]
async def double_critic_evaluate(
    self, question: str, sql: str, schema: str
) -> tuple[bool, str]:
    """Evaluate with double-critic to mitigate sycophancy."""
    
    # Query 1: Is this correct?
    correct_response = await self.generate(
        f"Is this SQL correct for the question?\n"
        f"Question: {question}\n"
        f"SQL: {sql}\n"
        f"Schema: {schema}\n"
        f"Answer YES or NO with explanation."
    )
    
    # Query 2: Is this incorrect?
    incorrect_response = await self.generate(
        f"Is this SQL INCORRECT for the question?\n"
        f"Question: {question}\n"
        f"SQL: {sql}\n"
        f"Schema: {schema}\n"
        f"Answer YES or NO with explanation."
    )
    
    # Accept only if: correct=YES AND incorrect=NO
    is_valid = (
        "YES" in correct_response.content.upper() and
        "NO" in incorrect_response.content.upper()
    )
    
    return is_valid, f"{correct_response.content}\n---\n{incorrect_response.content}"
\end{lstlisting}

\section{Output Format}

Generated examples are written to JSONL:

\begin{lstlisting}[language=Python,caption={JSONL output format}]
{"id": "sim-001", "question": "What is the total revenue...", 
 "sql": "SELECT SUM(amount) FROM...", "complexity_score": 0.3, ...}
{"id": "sim-002", "question": "Show all customers with...", 
 "sql": "SELECT * FROM customers WHERE...", "complexity_score": 0.7, ...}
\end{lstlisting}

% =============================================================================
% CHAPTER 6: ENVIRONMENT AND CLI
% =============================================================================
\newpage